% ============================================================
\subsection{Sliding Window}
% ============================================================

A sliding window is a contiguous span of a sequence, delimited by two indices
$\ell$ and $r$, together with a state summarizing the elements inside it. The
window advances by admitting the element at $r$ and evicting the element at
$\ell$. The technique applies when the state can be updated in $O(1)$ under
those two operations, so that recomputing it from scratch is never necessary.

The recognition cue is a request for a \emph{contiguous} subarray or substring
subject to a constraint. Contiguity is required. A problem about subsequences
or subsets admits reordering and belongs to a different paradigm.

% ------------------------------------------------------------
\subsubsection*{Why the cost is linear}

\begin{keybox}[The amortization argument]
Both indices are monotone. $r$ advances once per iteration of the outer loop
and $\ell$ never decreases, so each index takes at most $n$ values over the
entire run. The inner \texttt{while} loop is therefore bounded by $n$
increments \emph{in total} rather than $n$ increments per outer iteration.

The nesting is syntactic. The work is $O(n)$: every element is admitted exactly
once and evicted at most once. The naive alternative enumerates all $O(n^2)$
spans and pays the cost of the state on each.
\end{keybox}

The precondition is that $\ell$ may safely move forward. This holds when the
constraint is \textbf{monotone in the window}: if a window violates the
constraint, every window containing it also violates it. Under that property,
once $[\ell, r]$ fails, no window starting before $\ell$ and ending at $r$ or
later can succeed, so the discarded prefixes are provably useless.

% ------------------------------------------------------------
\subsubsection*{Fixed-size windows}

When the length $k$ is given, there is no contraction step. The window is
initialized once and then translated, one position at a time, by a paired
admit and evict.

\begin{lstlisting}[style=python, caption={Maximum sum over any window of length k}]
def max_window_sum(xs, k):
    if len(xs) < k:
        return None
    total = sum(xs[:k])           # the window is built once, in O(k)
    best  = total
    for r in range(k, len(xs)):
        total += xs[r] - xs[r - k]   # admit xs[r], evict xs[r-k]
        best = max(best, total)
    return best

max_window_sum([2, 1, 5, 1, 3, 2], 3)    # 9, from the span [5, 1, 3]
\end{lstlisting}

The single line \texttt{total += xs[r] - xs[r-k]} is the whole method. It
replaces an inner loop of length $k$, reducing $O(nk)$ to $O(n)$. The index
$r-k$ is the element leaving the window, since the window covering positions
$r-k+1$ through $r$ has just replaced the one covering $r-k$ through $r-1$.

\paragraph{A counter as the window state.} The state need not be a number.
Matching a pattern's permutations, described earlier under \texttt{Counter},
carries a frequency table as the state and compares it against a fixed target.

\begin{lstlisting}[style=python, caption={Every starting index at which a permutation of p occurs in t}]
from collections import Counter

def find_anagrams(t, p):
    k = len(p)
    need = Counter(p)
    have = Counter(t[:k])
    out  = [0] if have == need else []
    for r in range(k, len(t)):
        have[t[r]]     += 1              # admit
        have[t[r - k]] -= 1              # evict
        if have[t[r - k]] == 0:
            del have[t[r - k]]           # see the note below
        if have == need:
            out.append(r - k + 1)        # the window's left index
    return out

find_anagrams("cbaebabacd", "abc")       # [0, 6]
\end{lstlisting}

Each shift performs a constant amount of work and one dictionary comparison
over an alphabet of bounded size, giving $O(n)$.

\begin{keybox}[Zero-valued keys in a \texttt{Counter}]
Decrementing to zero leaves the key present with value \texttt{0}. Two
consequences follow, and they behave differently.

\medskip

\textbf{Equality.} On Python 3.10 and later, \texttt{Counter.\_\_eq\_\_}
compares counts over the union of the keys and treats a missing key as zero, so
\texttt{Counter(\{'a':1,'b':0\}) == Counter(\{'a':1\})} evaluates to
\texttt{True}. On Python 3.9 and earlier the class inherits
\texttt{dict.\_\_eq\_\_}, which compares key sets, and the same expression is
\texttt{False}.

\medskip

\textbf{Size.} \texttt{len(counter)} counts keys on every version. A key held at
zero still contributes to \texttt{len}, so any window that tests the number of
\emph{distinct} elements must delete the key on reaching zero.

\medskip

Deleting the exhausted key satisfies both, on every version.
\end{keybox}

% ------------------------------------------------------------
\subsubsection*{Variable-size windows}

When the length is unknown, $r$ expands the window on every iteration and an
inner loop moves $\ell$ forward. The two problem shapes differ in the test
governing that inner loop and in the point at which the answer is recorded.

\begin{center}
\begin{tabular}{lll}
\toprule
& \textbf{Longest valid window} & \textbf{Shortest sufficient window} \\
\midrule
Inner loop & \texttt{while} invariant is broken & \texttt{while} condition holds \\
Effect     & restore validity & test whether $\ell$ can be given up \\
Record at  & after the loop, on a valid window & inside the loop, before evicting \\
Answer     & \texttt{max(best, r - l + 1)}      & \texttt{min(best, r - l + 1)} \\
\bottomrule
\end{tabular}
\end{center}

\medskip

\paragraph{Longest valid window.} The window is repaired after each expansion,
so it is valid at the bottom of every outer iteration, which is where its length
is measured.

\begin{lstlisting}[style=python, caption={Longest substring with no repeated character}]
def longest_unique(s):
    seen = set()                  # the state: which characters are in the window
    left = best = 0
    for right, c in enumerate(s):
        while c in seen:          # admitting c would break the invariant
            seen.remove(s[left])  # evict from the left until it would not
            left += 1
        seen.add(c)               # admit
        best = max(best, right - left + 1)   # the window is valid here
    return best

longest_unique("abcabcbb")        # 3, from "abc"
longest_unique("pwwkew")          # 3, from "wke"
\end{lstlisting}

The eviction loop is written before the admission so that the invariant is never
violated inside the set. Admitting first and then repairing is equally correct
provided the repair condition is written to match, and the ordering must be
chosen once and held to.

\begin{lstlisting}[style=python, caption={Longest substring with at most k distinct characters}]
from collections import defaultdict

def longest_k_distinct(s, k):
    count = defaultdict(int)      # the state: a frequency table over the window
    left = best = 0
    for right, c in enumerate(s):
        count[c] += 1                        # admit
        while len(count) > k:                # invariant: at most k keys
            count[s[left]] -= 1
            if count[s[left]] == 0:
                del count[s[left]]           # required, or len(count) is wrong
            left += 1
        best = max(best, right - left + 1)
    return best

longest_k_distinct("eceba", 2)    # 3, from "ece"
\end{lstlisting}

Here the admission precedes the repair, because the invariant is stated over the
window's state rather than over the incoming element. The \texttt{del} is what
makes \texttt{len(count)} equal to the number of distinct characters currently
present.

\paragraph{Shortest sufficient window.} The inner loop contracts while the
target condition is still satisfied, which searches for the smallest left index
that keeps it satisfied. The length is recorded before each eviction, since the
window may cease to qualify immediately afterward.

\begin{lstlisting}[style=python, caption={Shortest subarray with sum at least t, over positive values}]
def min_length(xs, t):
    total = 0
    left  = 0
    best  = float('inf')
    for right, x in enumerate(xs):
        total += x                           # admit
        while total >= t:                    # the window already suffices
            best = min(best, right - left + 1)  # record, then try to shrink
            total -= xs[left]                # evict
            left  += 1
    return 0 if best == float('inf') else best

min_length([2, 3, 1, 2, 4, 3], 7)   # 2, from [4, 3]
min_length([1, 1, 1], 7)            # 0, no window qualifies
\end{lstlisting}

\texttt{float('inf')} as the initial value of a minimum removes the special case
for an empty candidate set. The sentinel is then translated to whatever the
problem specifies for ``no such window,'' which is commonly \texttt{0}.

% ------------------------------------------------------------
\subsubsection*{When the window does not apply}

\begin{keybox}[The monotonicity precondition]
\texttt{min\_length} above is correct only because every element is positive.
Positivity makes \texttt{total} nondecreasing in the window's length, so
\texttt{total >= t} is monotone and a failing prefix can be discarded
permanently.

Introduce a negative value and that reasoning collapses. Extending a window may
now decrease its sum, and a left index abandoned earlier could belong to the
optimal answer. On \texttt{[3, -2, 5]} with $t = 6$, the window
\texttt{[3, -2, 5]} qualifies while neither the prefix \texttt{[3, -2]} nor the
suffix \texttt{[5]} taken alone reveals it by contraction.

The replacement is prefix sums with a hash map, which considers every left
endpoint rather than only the ones a monotone rule preserves. The check for a
negative value in the constraints is therefore the first thing to perform on any
subarray-sum problem.
\end{keybox}

% ------------------------------------------------------------
\subsubsection*{Common errors}

\begin{itemize}
  \item The window length is \texttt{right - left + 1}. Omitting the
        \texttt{+1} is the standard off-by-one, since both endpoints are
        inclusive.
  \item Recording the answer at the wrong point. A longest-window problem
        measures after the repair loop. A shortest-window problem measures
        inside the contraction loop, before the eviction.
  \item An \texttt{if} where a \texttt{while} is required. One eviction restores
        the invariant only when a single admission can break it by exactly one
        unit. The general repair is a loop.
  \item Failing to evict from the state. Advancing \texttt{left} without
        removing \texttt{xs[left]} from the set, counter, or running sum leaves
        the state describing a window that no longer exists.
  \item Retaining zero-valued keys in a \texttt{Counter} or
        \texttt{defaultdict} while testing \texttt{len} against a distinct-count
        bound.
  \item Recomputing the state over the span, as in \texttt{sum(xs[left:right+1])}
        inside the loop. This restores the $O(n^2)$ cost the window exists to
        remove.
\end{itemize}
